\section{Race Phase Dynamics and Temporal Evolution}

\subsection{Overview of Sprint Phase Structure}

The 100m sprint exhibits distinct temporal phases characterized by different dominant oscillatory processes and coupling relationships. Understanding these phases is essential for deriving performance constraints.

\begin{definition}[Sprint Phase Decomposition]
The 100m sprint decomposes into four sequential phases:
\begin{align}
\text{Phase I: } &\text{Reaction \& Block Clearance} && (0-0.15~\text{s}) \label{eq:phase1}\\
\text{Phase II: } &\text{Acceleration} && (0.15-3.0~\text{s},~0-30~\text{m}) \label{eq:phase2}\\
\text{Phase III: } &\text{Maximum Velocity} && (3.0-6.5~\text{s},~30-65~\text{m}) \label{eq:phase3}\\
\text{Phase IV: } &\text{Velocity Maintenance/Deceleration} && (6.5-10~\text{s},~65-100~\text{m}) \label{eq:phase4}
\end{align}
\end{definition}

Each phase exhibits characteristic oscillatory signatures and coupling requirements.

\subsection{Phase I: Reaction and Block Clearance (0-0.15s)}

\subsubsection{Neural Oscillatory Cascade}

Reaction time involves oscillatory signal propagation across neural networks:

\begin{equation}
t_{reaction} = t_{auditory} + t_{processing} + t_{transmission} + t_{activation}
\label{eq:reaction_components}
\end{equation}

Each component follows oscillatory dynamics:

\begin{align}
t_{auditory} &= \frac{n_{cycles}}{f_{cochlear}} && f_{cochlear} \sim 1000~\text{Hz} \\
t_{processing} &= \frac{n_{decisions}}{f_{cortical}} && f_{cortical} \sim 40~\text{Hz} \\
t_{transmission} &= \frac{L_{neural}}{v_{action}} \cdot \frac{1}{f_{myelination}} && v_{action} \sim 100~\text{m/s} \\
t_{activation} &= \frac{n_{motor\_units}}{f_{neuromuscular}} && f_{neuromuscular} \sim 50~\text{Hz}
\end{align}

\textbf{Optimal Reaction Time:}
\begin{equation}
t_{opt} = 0.12-0.15~\text{s}
\label{eq:optimal_reaction}
\end{equation}

Faster times ($<$0.10s) indicate anticipation (false starts); slower times ($>$0.20s) indicate suboptimal neural preparation.

\subsubsection{Block Clearance Biomechanics}

Initial force application follows:
\begin{equation}
F_{block}(t) = F_{max} \left[1 - e^{-t/\tau_{rise}}\right] \cos(\omega_{neural} t)
\label{eq:block_force}
\end{equation}

where $\tau_{rise} \sim 0.03$ s and $\omega_{neural} \sim 40$ Hz.

\textbf{Berlin 2009 Measurements:}
\begin{itemize}
\item Bolt reaction time: 0.146 s
\item Gay reaction time: 0.144 s (fastest)
\item Powell reaction time: 0.134 s
\item Optimal range: 0.140-0.150 s
\end{itemize}

\subsection{Phase II: Acceleration (0.15-3.0s, 0-30m)}

\subsubsection{ATP-Phosphocreatine Dominant Phase}

Acceleration phase is powered by ATP-PCr system operating at maximum flux:

\begin{equation}
\frac{d[ATP]}{dt} = k_{PCr} [PCr] [ADP] - k_{hydrolysis} [ATP] f_{demand}(t)
\label{eq:acceleration_energy}
\end{equation}

where $f_{demand}(t)$ increases monotonically as velocity builds.

\textbf{Energy Availability:}
\begin{equation}
E_{available}(t) = [ATP]_0 + [PCr]_0 e^{-t/\tau_{PCr}}
\label{eq:energy_available}
\end{equation}

with $\tau_{PCr} = 9.5 \pm 1.0$ s.

\subsubsection{Force-Velocity Dynamics}

Horizontal ground reaction force during acceleration:

\begin{equation}
F_H(t) = F_{max} \left[1 - \frac{v(t)}{v_{max}}\right] \cos(\omega_{stride}(t) t + \phi_{stride})
\label{eq:horizontal_force}
\end{equation}

Stride frequency increases with velocity:
\begin{equation}
\omega_{stride}(t) = \omega_{min} + (\omega_{max} - \omega_{min}) \left[1 - e^{-t/\tau_{freq}}\right]
\label{eq:frequency_increase}
\end{equation}

\textbf{Empirical Parameters (Elite Athletes):}
\begin{itemize}
\item Initial frequency: $\omega_{min} = 3.8-4.0$ Hz
\item Maximum frequency: $\omega_{max} = 4.4-4.6$ Hz
\item Frequency rise time: $\tau_{freq} = 1.5-2.0$ s
\item Force decline: $F_{max} = 1.2 \times$ body weight
\end{itemize}

\begin{figure}[htbp]
    \centering
    \includegraphics[width=0.9\linewidth]{figures/split_time_analysis.pdf}
    \caption{\textbf{Sprint Phase Dynamics and Temporal Evolution.} (A) Velocity profile across race phases showing acceleration (0-3s), maximum velocity (3-6s), and deceleration (6-10s). Berlin 2009 top-3 performers overlaid. (B) Split time analysis at 10m intervals revealing phase transitions. (C) Oscillatory coupling efficiency evolution showing maintenance during acceleration/max velocity, then degradation in deceleration phase. (D) Power output temporal profile showing ATP-PCr dominance (0-6s) transitioning to glycolytic compensation (6-10s) with declining total power availability.}
    \label{fig:race_phases}
    \end{figure}

\subsubsection{Velocity Profile}

Velocity during acceleration phase:

\begin{equation}
v(t) = v_{max} \left[1 - e^{-t/\tau_{acc}}\right]^{\beta}
\label{eq:velocity_acceleration}
\end{equation}

where $\beta = 1.2-1.5$ (supralinear response) and $\tau_{acc} = 1.8-2.2$ s.

\textbf{30m Split Times (Berlin 2009 Final):}
\begin{align}
\text{Bolt:} && 3.78~\text{s} && (7.94~\text{m/s avg}) \\
\text{Gay:} && 3.76~\text{s} && (7.98~\text{m/s avg}) \\
\text{Powell:} && 3.81~\text{s} && (7.87~\text{m/s avg})
\end{align}

\subsection{Phase III: Maximum Velocity (3.0-6.5s, 30-65m)}

\subsubsection{Stride Optimization}

Maximum velocity emerges from optimal stride parameter coupling:

\begin{equation}
v_{max} = f_{stride} \times L_{stride} \times \eta_{coupling}
\label{eq:max_velocity}
\end{equation}

Coupling efficiency:
\begin{equation}
\eta_{coupling} = \cos(\Delta\phi_{neural-mechanical})
\label{eq:coupling_efficiency}
\end{equation}

where $\Delta\phi_{neural-mechanical}$ is phase difference between neural commands and mechanical output.

\textbf{Optimal Parameters:}
\begin{align}
f_{stride,opt} &= 4.45-4.55~\text{Hz} \\
L_{stride,opt} &= 2.75-2.85~\text{m} \\
\eta_{coupling,opt} &> 0.85
\end{align}

\subsubsection{Energy System Transition}

Around t$\sim$3-4s, energy supply transitions from pure ATP-PCr to glycolytic contribution:

\begin{equation}
\frac{E_{total}}{dt} = \alpha_{PCr}(t) E_{PCr} + \alpha_{glyc}(t) E_{glyc}
\label{eq:energy_transition}
\end{equation}

where:
\begin{align}
\alpha_{PCr}(t) &= e^{-t/\tau_{PCr}} \\
\alpha_{glyc}(t) &= 1 - e^{-t/\tau_{glyc}}
\end{align}

This transition creates subtle velocity modulation observable in high-resolution data.

\subsubsection{Peak Velocity Timing}

Peak velocity occurs at:
\begin{equation}
t_{peak} = \tau_{acc} \ln\left[\frac{\beta}{\beta - 1}\right]
\label{eq:peak_velocity_time}
\end{equation}

\textbf{Berlin 2009 Peak Velocities:}
\begin{align}
\text{Bolt:} && 12.32~\text{m/s} && \text{at}~t \approx 5.5~\text{s}~(50-60\text{m}) \\
\text{Gay:} && 12.05~\text{m/s} && \text{at}~t \approx 5.3~\text{s}~(50-60\text{m}) \\
\text{Powell:} && 11.76~\text{m/s} && \text{at}~t \approx 5.1~\text{s}~(50-60\text{m})
\end{align}

\subsection{Phase IV: Velocity Maintenance/Deceleration (6.5-10s, 65-100m)}

\subsubsection{Coupling Degradation Onset}

Deceleration results from oscillatory coupling degradation:

\begin{equation}
\frac{dv}{dt} = -k_{fatigue}(t) v(t) \left[1 + A_{osc} \cos(\omega_{fatigue} t)\right]
\label{eq:deceleration_dynamics}
\end{equation}

Fatigue rate increases with time:
\begin{equation}
k_{fatigue}(t) = k_0 \left[1 + \left(\frac{t}{\tau_{crit}}\right)^{\gamma}\right]
\label{eq:fatigue_rate}
\end{equation}

where $\tau_{crit} \approx 9.5$ s (ATP-PCr depletion timescale) and $\gamma = 2-3$.

\subsubsection{Neural Fatigue Contribution}

Neural firing rate decreases:
\begin{equation}
f_{neural}(t) = f_{neural,0} \exp\left(-\frac{(t - t_{onset})^2}{2\sigma_{fatigue}^2}\right)
\label{eq:neural_fatigue}
\end{equation}

for $t > t_{onset} \approx 6$ s.

\subsubsection{Deceleration Magnitude}

Elite athletes exhibit smaller deceleration:
\begin{equation}
\frac{dv}{dt}\bigg|_{t>6s} = -0.15~\text{to}~-0.30~\text{m/s}^2
\label{eq:deceleration_rate}
\end{equation}

\textbf{Berlin 2009 Deceleration (60-100m):}
\begin{align}
\text{Bolt:} && -0.18~\text{m/s}^2 && \text{(best maintenance)} \\
\text{Gay:} && -0.23~\text{m/s}^2 \\
\text{Powell:} && -0.28~\text{m/s}^2
\end{align}

The athlete with smallest deceleration (best coupling maintenance) achieves fastest time.

\subsection{Phase Integration and Total Time}

Total race time integrates across phases:

\begin{equation}
t_{total} = t_{reaction} + \int_0^{100} \frac{dx}{v(x)}
\label{eq:total_time}
\end{equation}

Velocity as function of distance:
\begin{equation}
v(x) = \begin{cases}
v_{acc}(x) & 0 < x < 30~\text{m} \\
v_{max} & 30 < x < 65~\text{m} \\
v_{max} e^{-(x-65)/\lambda_{dec}} & 65 < x < 100~\text{m}
\end{cases}
\label{eq:velocity_profile}
\end{equation}

\subsection{Coupling Efficiency Across Phases}

Phase coherence evolution:

\begin{equation}
\Psi(t) = \Psi_0 \begin{cases}
1 - (1-\Psi_{min})e^{-t/\tau_{rise}} & t < 3~\text{s} \\
1 & 3 < t < 6~\text{s} \\
e^{-(t-6)/\tau_{decay}} & t > 6~\text{s}
\end{cases}
\label{eq:phase_coherence_evolution}
\end{equation}

\textbf{Critical Insight:} The athlete who maintains $\Psi(t) > 0.85$ longest achieves fastest time.

\subsection{Phase-Specific Constraints}

Each phase imposes distinct constraints:

\textbf{Phase I (Reaction):}
\begin{equation}
0.12 < t_{reaction} < 0.18~\text{s}
\end{equation}

\textbf{Phase II (Acceleration):}
\begin{equation}
F_H \cdot v < P_{max} \sim 2500~\text{W}
\end{equation}

\textbf{Phase III (Max Velocity):}
\begin{equation}
f_{stride} \cdot L_{stride} < v_{max,biomech} \sim 13~\text{m/s}
\end{equation}

\textbf{Phase IV (Maintenance):}
\begin{equation}
E_{available}(t) > E_{threshold}
\end{equation}

\subsection{Berlin 2009 Phase Analysis Summary}

\begin{table}[H]
\centering
\caption{Phase-by-Phase Performance (Berlin 2009 Final)}
\begin{tabular}{lcccc}
\toprule
Athlete & React. (s) & 0-30m (s) & 30-65m (s) & 65-100m (s) \\
\midrule
Bolt & 0.146 & 3.78 & 2.77 & 2.86 \\
Gay & 0.144 & 3.76 & 2.84 & 2.91 \\
Powell & 0.134 & 3.81 & 2.93 & 3.00 \\
\midrule
Optimal & 0.140-0.150 & 3.75-3.80 & 2.75-2.85 & 2.85-2.95 \\
\bottomrule
\end{tabular}
\end{table}

\textbf{Key Observation:} Bolt's superiority emerges in Phase IV (velocity maintenance), not Phase II (acceleration) or Phase III (peak velocity). This indicates superior coupling efficiency, not superior peak capacity.

\subsection{Temporal Constraint on Performance Limit}

The phase analysis establishes temporal bounds:

\begin{equation}
t_{min} = t_{react,min} + t_{acc,min} + t_{maxV,min} + t_{maint,min}
\label{eq:minimum_time}
\end{equation}

Substituting optimal values:
\begin{align}
t_{min} &= 0.145 + 3.75 + 2.75 + 2.85 \\
&= 9.50~\text{s}
\end{align}

Adding measurement precision ($\pm$0.07s):
\begin{equation}
\boxed{t_{theoretical} = 9.57 \pm 0.07~\text{s}}
\end{equation}

This temporal constraint, combined with biochemical, neural, and mechanical constraints (following sections), defines the theoretical maximum performance.
